\documentclass[11pt,a4paper,sans]{moderncv}

\moderncvstyle{classic}
\moderncvcolor{blue}
\usepackage[scale=0.84]{geometry}
\usepackage{fontspec}
\usepackage{xeCJK}
\usepackage{enumitem}
\usepackage{xurl}

\setmainfont{Times New Roman}
\setsansfont{Helvetica}
\setmonofont{Menlo}
\setCJKmainfont{Songti SC}
\setCJKsansfont{Noto Sans CJK SC}

\setlist[itemize]{leftmargin=*, itemsep=0.18em, topsep=0.25em}
\setlength{\hintscolumnwidth}{3.1cm}
\renewcommand*{\namefont}{\fontsize{28}{32}\mdseries\upshape}
\recomputelengths

\name{Weiwei}{Chen}
\address{Institute of Logic and Cognition, Department of Philosophy}{Sun Yat-sen University, Guangzhou 510275, China}
\email{chenweiwei@mail.sysu.edu.cn}
\homepage{chenww.com}
\extrainfo{ORCID: \href{https://orcid.org/0000-0002-3447-8163}{0000-0002-3447-8163} \quad
Google Scholar: \href{https://scholar.google.com/citations?user=I8ISBZAAAAAJ}{I8ISBZAAAAAJ} \quad
DBLP: \href{https://dblp.org/pid/68/925-6.html}{Weiwei Chen 0006}}

\begin{document}
\makecvtitle

\section{Research Interests}
\cvitem{}{Formal argumentation; collective argumentation; social choice theory; logic and AI; computational models of argument; dynamic and strategic collective decision-making.}

\section{Current Appointment}
\cventry{2024--present}{Associate Professor and MA Supervisor}{Institute of Logic and Cognition, Department of Philosophy, Sun Yat-sen University}{Guangzhou, China}{}{}

\section{Employment}
\cventry{2019--2024}{Postdoctoral Researcher; later Associate Researcher}{Institute of Logic and Cognition, Department of Philosophy, Sun Yat-sen University}{Guangzhou, China}{}{}

\section{Education}
\cventry{2015--2019}{PhD in Logic}{Institute of Logic and Cognition, Department of Philosophy, Sun Yat-sen University}{Guangzhou, China}{}{}
\cventry{2016--2017}{Visiting PhD Student}{Institute for Logic, Language and Computation, University of Amsterdam}{Amsterdam, Netherlands}{}{}
\cventry{2013--2015}{MSc in Logic}{Institute of Logic and Cognition, Department of Philosophy, Sun Yat-sen University}{Guangzhou, China}{}{}
\cventry{2006--2010}{BA in Software Engineering}{Guangzhou University}{Guangzhou, China}{}{}

\section{Selected Publications}
\cvitem{2019}{Weiwei Chen and Ulle Endriss. \href{https://doi.org/10.1016/j.artint.2018.10.006}{Preservation of Semantic Properties in Collective Argumentation: The Case of Aggregating Abstract Argumentation Frameworks}. \emph{Artificial Intelligence}, 269: 27--48.}
\cvitem{2021}{Weiwei Chen. \href{https://doi.org/10.1093/logcom/exab011}{Guaranteeing Admissibility of Abstract Argumentation Frameworks with Rationality and Feasibility Constraints}. \emph{Journal of Logic and Computation}, 31(8): 2133--2158.}
\cvitem{2023}{Weiwei Chen. \href{https://doi.org/10.1093/logcom/exac100}{Collective Argumentation with Topological Restrictions: The Case of Aggregating Abstract Argumentation Frameworks}. \emph{Journal of Logic and Computation}, 33(2): 319--343.}
\cvitem{2024}{Weiwei Chen and Shier Ju. \href{https://doi.org/10.1016/j.ijar.2024.109234}{Dynamic Collective Argumentation: Constructing the Revision and Contraction Operators}. \emph{International Journal of Approximate Reasoning}, 172: 109234.}
\cvitem{2026}{Xiaolong Liu, Weiwei Chen, and Sylvie Doutre. \href{https://doi.org/10.1093/logcom/exag003}{Argumentation Semantics Based on Solid Admissibility and Preservation of Semantic Properties}. \emph{Journal of Logic and Computation}, 36(2): exag003.}

\section{Books and Book Chapters}
\cvitem{Forthcoming}{陈伟伟. 《论辩聚合：社会选择理论的新视角》. 人民出版社.}
\cvitem{2022}{刘小龙、陈伟伟. 《论证挖掘与论证形式化》第十二章. 鞠实儿等著, 科学出版社.}
\cvitem{2020}{陈伟伟 and Ulle Endriss. 《融合与修正：跨文化交流的逻辑与认知研究》第五章. 鞠实儿等著, 经济科学出版社.}

\section{Journal Articles}
\cvitem{2026}{Xiaolong Liu, Weiwei Chen, and Sylvie Doutre. \href{https://doi.org/10.1093/logcom/exag003}{Argumentation Semantics Based on Solid Admissibility and Preservation of Semantic Properties}. \emph{Journal of Logic and Computation}, 36(2): exag003.}
\cvitem{2024}{Weiwei Chen and Shier Ju. \href{https://doi.org/10.1016/j.ijar.2024.109234}{Dynamic Collective Argumentation: Constructing the Revision and Contraction Operators}. \emph{International Journal of Approximate Reasoning}, 172: 109234.}
\cvitem{2024}{Zhe Yu, Shier Ju, and Weiwei Chen. \href{https://doi.org/10.1093/logcom/exad063}{Context-based Argumentation Frameworks and Multi-Agent Consensus Building}. \emph{Journal of Logic and Computation}, 34(2): 199--228.}
\cvitem{2023}{Weiwei Chen. \href{https://doi.org/10.1093/logcom/exac100}{Collective Argumentation with Topological Restrictions: The Case of Aggregating Abstract Argumentation Frameworks}. \emph{Journal of Logic and Computation}, 33(2): 319--343.}
\cvitem{2021}{Weiwei Chen. \href{https://doi.org/10.1093/logcom/exab011}{Guaranteeing Admissibility of Abstract Argumentation Frameworks with Rationality and Feasibility Constraints}. \emph{Journal of Logic and Computation}, 31(8): 2133--2158.}
\cvitem{2019}{Weiwei Chen and Ulle Endriss. \href{https://doi.org/10.1016/j.artint.2018.10.006}{Preservation of Semantic Properties in Collective Argumentation: The Case of Aggregating Abstract Argumentation Frameworks}. \emph{Artificial Intelligence}, 269: 27--48.}

\section{Conference and Workshop Papers}
\cvitem{2025}{Xiaolong Liu and Weiwei Chen. \href{https://doi.org/10.1007/978-981-96-7956-0_16}{Proportional Acceptability of Arguments}. In \emph{Logic and Argumentation, CLAR 2025}, LNCS 15712, Springer, pp. 254--271.}
\cvitem{2024}{Xiaolong Liu and Weiwei Chen. On the Proportional Acceptability of Arguments in Abstract Argumentation. In \emph{Proceedings of the 10th International Conference on Computational Models of Argument}, IOS Press.}
\cvitem{2021}{Weiwei Chen. \href{https://doi.org/10.1007/978-3-030-89391-0_5}{Collective Argumentation with Topological Restrictions}. In \emph{Logic and Argumentation, CLAR 2021}, LNCS 13040, Springer, pp. 79--93.}
\cvitem{2021}{Xiaolong Liu and Weiwei Chen. \href{https://doi.org/10.1007/978-3-030-82254-5_11}{Solid Semantics for Abstract Argumentation Frameworks and the Preservation of Solid Semantic Properties}. In \emph{EUMAS 2021}, LNCS 12802, Springer, pp. 178--193.}
\cvitem{2021}{Weiwei Chen. \href{https://doi.org/10.4204/EPTCS.335.9}{Collective Argumentation: The Case of Aggregating Support-Relations of Bipolar Argumentation Frameworks}. \emph{Electronic Proceedings in Theoretical Computer Science}, 335: 87--102.}
\cvitem{2021}{Xiaolong Liu and Weiwei Chen. Solid Semantics and Extension Aggregation Using Quota Rules under Integrity Constraints. In \emph{Proceedings of the 20th International Conference on Autonomous Agents and Multiagent Systems (AAMAS 2021)}.}
\cvitem{2020}{Weiwei Chen. Aggregation of Support-Relations of Bipolar Argumentation Frameworks. In \emph{Proceedings of the 19th International Conference on Autonomous Agents and Multiagent Systems (AAMAS 2020)}.}
\cvitem{2020}{Weiwei Chen. \href{https://doi.org/10.1007/978-3-030-44638-3_15}{Preservation of Admissibility with Rationality and Feasibility Constraints}. In \emph{Logic and Argumentation, CLAR 2020}, LNAI 12061, Springer, pp. 245--258.}
\cvitem{2018}{Weiwei Chen and Ulle Endriss. \href{https://doi.org/10.3233/978-1-61499-906-5-425}{Aggregating Alternative Extensions of Abstract Argumentation Frameworks: Preservation Results for Quota Rules}. In \emph{Computational Models of Argument, COMMA 2018}, Frontiers in Artificial Intelligence and Applications 305, IOS Press, pp. 425--436.}
\cvitem{2017}{Weiwei Chen and Ulle Endriss. \href{https://doi.org/10.4204/EPTCS.251.9}{Preservation of Semantic Properties during the Aggregation of Abstract Argumentation Frameworks}. \emph{Electronic Proceedings in Theoretical Computer Science}, 251: 118--133.}

\section{Research Grants}
\cvitem{2024}{Principal Investigator, 公理化群体论辩的可能性和防操纵性研究. Natural Science Foundation of Guangdong, China, 2024A1515012540.}
\cvitem{2022}{Principal Investigator, 基于社会选择方法的论辩聚合理论研究. National Social Science Fund of China, 22FZXB093.}
\cvitem{2022}{Principal Investigator, 公理化群体论辩理论. Ministry of Education of China, Humanities and Social Sciences Youth Fund, 22YJC72040001.}
\cvitem{2019}{Principal Investigator, 形式论辩中的双极论辩框架与分级论辩语义聚合研究. China Postdoctoral Science Foundation, 2019M663352.}

\section{Selected Talks}
\cvitem{2024}{Constructing the Revision and Contraction Operators for Dynamic Collective Argumentation. The 3rd International Workshop on Logic for Multi-Agent Systems, Zhejiang University, June 2024.}
\cvitem{2024}{Dynamically Rational Collective Argumentation. The 3rd International Workshop on Logic and Philosophy, Tsinghua University, May 2024.}
\cvitem{2024}{Collective Argumentation: Static and Dynamic Perspectives. Seminar at PKU Center of Logic, Language, and Cognition, Peking University, March 2024.}
\cvitem{2023}{Dynamic Collective Argumentation. Trends in Logic, Toruń, Poland, November 2023.}
\cvitem{2022}{Applying the Axiomatic Method to Collective Argumentation. Individual and Collective Reasoning Group Seminar, University of Luxembourg, February 2022.}
\cvitem{2021}{Collective Argumentation with Topological Restrictions. CLAR 2021, online, October 2021.}
\cvitem{2021}{Collective Argumentation: The Case of Aggregation of Support-Relations of Bipolar Argumentation Frameworks. TARK 2021, online, June 2021.}
\cvitem{2020}{Preservation of Admissibility with Rationality and Feasibility Constraints. LILaC Team Seminar, IRIT, France, online, December 2020.}
\cvitem{2020}{Aggregation of Support-Relations of Bipolar Argumentation Frameworks. AAMAS 2020, online, May 2020.}
\cvitem{2018}{Aggregating Extensions of Abstract Argumentation Frameworks. Computational Social Choice Seminar, ILLC, University of Amsterdam, September 2018.}
\cvitem{2017}{What Can We Expect from Consensus Decision-making? Cool Logic Seminar, University of Amsterdam, September 2017.}
\cvitem{2017}{Aggregation of Abstract Argumentation Frameworks. TARK 2017, Liverpool, United Kingdom, July 2017.}

\section{Professional Service}
\cvitem{Program committee}{ECAI 2024, 2025; LNGAI 2023, 2024; AAMAS 2023, 2024, 2025; PRICAI 2023, 2024; CLAR 2020, 2023, 2025; IJCAI 2020, 2025.}
\cvitem{Reviewer}{\emph{Autonomous Agents and Multi-Agent Systems}; \emph{Journal of Applied Non-Classical Logics}; \emph{Studies in Logic} (逻辑学研究).}

\section{Academic Profiles}
\cvitem{Homepage}{\href{https://chenww.com/}{https://chenww.com/}}
\cvitem{Faculty page}{\href{https://philosophy.sysu.edu.cn/teacher/ChenWeiwei}{https://philosophy.sysu.edu.cn/teacher/ChenWeiwei}}
\cvitem{Google Scholar}{\href{https://scholar.google.com/citations?user=I8ISBZAAAAAJ&hl=en}{https://scholar.google.com/citations?user=I8ISBZAAAAAJ}}
\cvitem{ORCID}{\href{https://orcid.org/0000-0002-3447-8163}{0000-0002-3447-8163}}
\cvitem{DBLP}{\href{https://dblp.org/pid/68/925-6.html}{Weiwei Chen 0006} (partial index of computer-science-related publications)}

\vfill
\begin{center}
{\small Last updated: May 19, 2026}
\end{center}

\end{document}
